% Generated draft: requires mathematical and source review.
\documentclass[11pt]{article}
\usepackage[margin=25mm]{geometry}
\usepackage{fontspec}
\setmainfont{Latin Modern Roman}
\usepackage{amsmath,amssymb,mathtools,unicode-math}
\setmathfont{Latin Modern Math}
\usepackage{enumitem,needspace,xurl,hyperref}
\setlength{\parindent}{0pt}
\setlength{\parskip}{5pt}
\setlength{\emergencystretch}{3em}
\newcommand{\N}{\mathbb{N}}
\newcommand{\Z}{\mathbb{Z}}
\newcommand{\Q}{\mathbb{Q}}
\newcommand{\R}{\mathbb{R}}
\newcommand{\C}{\mathbb{C}}
\newcommand{\F}{\mathbb{F}}
\newcommand{\A}{\mathbb{A}}
\newcommand{\PP}{\mathbb{P}}
\newcommand{\E}{\mathbb{E}}
\newcommand{\eps}{\varepsilon}
\newcommand{\dd}{\,\mathrm{d}}
\begin{document}
\section*{UnsolvedMath ID: 10600018}
\textbf{Source problem code:} catalogue entry
\subsection{Definitions and mathematical statement}
A tame virtual link may conventionally be represented by a finite virtual link diagram, considered modulo the classical Reidemeister moves and the virtual detour moves. Write $L_1\sim_{\mathrm t}L_2$ for this equivalence. The source does not specify an orientation convention or a model of wild virtual links.

Problem-list item 18 asks to construct a category of wild virtual knots and establish its axioms. This requires specifying its objects, morphisms, equivalence relation, and the relationship of its tame objects to ordinary tame virtual knots and links. In particular, seek a theorem specifying when a wild equivalence between tame virtual links entails a tame equivalence. If $\iota$ denotes the proposed passage to the wild theory and $\sim_{\mathrm w}$ its equivalence, the requested implication has the form
\[
 \iota(L_1)\sim_{\mathrm w}\iota(L_2)
 \quad\Longrightarrow\quad L_1\sim_{\mathrm t}L_2
\]
for tame $L_1,L_2$, under hypotheses on the wild theory and its equivalences that must be stated as part of the construction. The source asks when this implication holds; it does not supply a universal assertion for an already defined wild category. In particular, no topology on infinite diagrams, ambient category, admissible accumulation behavior, or wild equivalence moves are provided.
\subsection{Short English statement}
Define a theory and category of wild virtual knots, then determine precise conditions under which a wild equivalence between tame virtual links implies their usual tame equivalence.
\subsection{Sources}
\begin{itemize}
\item Local registry: \url{data/unsolvedmath/problems.json}, permanent UnsolvedMath ID $10600018$, exact source code \texttt{catalogue entry}, ``Virtual-knot problem 18 --- Wild Virtuals.''
\item \emph{Unsolved Problems in Virtual Knot Theory and Combinatorial Knot Theory}, revised survey, problem-list item 18: \url{https://arxiv.org/abs/1409.2823}. Supplied access date: 2026-07-29; extraction: source-tex. The problem section contains 60 top-level items and retains compound items intact.
\item No additional primary reference or author attribution is supplied for this problem. The registry's statement that its searches found no relevant literature is an unverified search assessment.
\end{itemize}
\end{document}
